% TOP_PROBLEM: {"id":30007690,"problem_number":"LOCAL-30007690","title":"Privacy regulation welfare","category_id":21,"category":{"id":21,"name":"economics","display_name":"Economics","description":"Open economic theory statements, published research agendas, and proposed scoped specifications; question provenance and historical priority are qualified per record.","slug":"economics","order_index":21,"created_at":"2026-10-08T00:00:00Z"},"status":"research_question","external_url":null,"legacy_ids":[],"published":false,"statement":"Estimate dynamic welfare effects of a specified privacy-enforcement change, including privacy benefits, compliance costs, entry and innovation.","economics":{"original_id":179,"field":"Digital, platform and data economics","kind":"empirical","formalization_basis":"adapted_research_question","source_status":"research_agenda","priority_status":"earliest_found","priority_note":"Earliest relevant explicit question or agenda passage located in this audit. No claim of first-ever formulation; earlier versions and later solutions were not exhaustively traced.","earliest_verified_reference":{"authors":["Garrett A. Johnson"],"title":"Economic Research on Privacy Regulation: Lessons from the GDPR and Beyond","year":2024,"url":"https://gwern.net/doc/cs/security/2024-johnson.pdf","doi":null,"locator":"Section 4.6, printed pp. 120–121, intended consumer benefits and regulatory design","evidence_summary":"Requests rigorous privacy-benefit measurement and study of effective regulation and unintended effects."},"current_reference":{"authors":["Garrett A. Johnson"],"title":"Economic Research on Privacy Regulation: Lessons from the GDPR and Beyond","year":2024,"url":"https://gwern.net/doc/cs/security/2024-johnson.pdf","doi":null,"locator":"Section 4.6, printed pp. 120–121, intended consumer benefits and regulatory design","evidence_summary":"Requests rigorous privacy-benefit measurement and study of effective regulation and unintended effects."},"formalization_note":"The source verifies a broader question or research agenda; the population, estimand, equations and restrictions here are our scoped adaptation. This is not a quoted canonical conjecture, and the audit does not prove contemporary unsolved status for every special case."},"statusQualification":"Published question or research agenda verified at the cited locator. The mathematical specification is adapted for this catalogue and includes additional assumptions; source publication does not establish first-ever priority or certify this exact model remains unsolved. The source verifies a broader question or research agenda; the population, estimand, equations and restrictions here are our scoped adaptation. This is not a quoted canonical conjecture, and the audit does not prove contemporary unsolved status for every special case.","statusReviewedAt":"2026-10-08"}
% FIRST_POSTED: 2026-10-08
% ENTRY_KIND: statement_only
% !TeX program = lualatex
\documentclass[11pt]{article}
\usepackage{fontspec}
\usepackage[a4paper,margin=25mm]{geometry}
\usepackage{amsmath,amssymb,mathtools}
\usepackage{xurl}
\usepackage[hidelinks]{hyperref}
\newcommand{\sref}[2]{\hyperlink{source:#1:#2}{[S#2]}}
\setlength{\emergencystretch}{3em}
\begin{document}
% problemId:30007690
\section*{Privacy regulation welfare}\label{30007690}
\subsection{Definitions and mathematical statement}
Fix European Union consumers and firms providing online retail services during a five-year period after a specified privacy-enforcement change. Define two regulatory regimes \(r\in\{0,1\}\), differing only in that enforcement policy, with their cross-border scope specified. Under either regime firms can alter data use, enter, exit and innovate. Let \(B_i(r)\) be consumer \(i\)'s monetized privacy and service benefit net of payments; let \(\Pi(r)\) be aggregate firm profit net of production, compliance and innovation costs; and let \(K(r)\) be real public enforcement cost. Define \(W(r)=\sum_i B_i(r)+\Pi(r)-K(r)\), summing over the fixed eligible-consumer population and including consumers of exiting firms. The target is \(\Delta W=W(1)-W(0)\), decomposed into privacy, service, market-structure and real compliance-cost components using a stated counterfactual ordering. Identify privacy benefits through audited data flows and informed valuation, while accounting for consent-related missing observations. Establish credible control jurisdictions or partial-identification restrictions when firms change behavior globally. Report the uncertainty associated with compliance, enforcement intensity and unobserved consumer harms. This is a dynamic welfare agenda for a specific policy change; it neither treats GDPR adoption as a clean universal experiment nor counts reduced measured tracking as a complete welfare assessment.

\par\textbf{Formulation and scope.} The source verifies a broader question or research agenda; the population, estimand, equations and restrictions here are our scoped adaptation. This is not a quoted canonical conjecture, and the audit does not prove contemporary unsolved status for every special case.

\par\textbf{Open-question provenance.} An explicit question or research agenda was verified at the source locator. This does not by itself certify that every later version remains unresolved \sref{30007690}{1}.

\par\textbf{Historical priority.} Earliest relevant explicit question or agenda passage located in this audit. No claim of first-ever formulation; earlier versions and later solutions were not exhaustively traced.

\subsection{Short English statement}
Estimate dynamic welfare effects of a specified privacy-enforcement change, including privacy benefits, compliance costs, entry and innovation.

\subsection{Sources}
\begin{itemize}\raggedright
\item[\textnormal{[S1]}]\hypertarget{source:30007690:1}{} Garrett A. Johnson. 2024. Economic Research on Privacy Regulation: Lessons from the GDPR and Beyond. Section 4.6, printed pp. 120–121, intended consumer benefits and regulatory design.
\url{https://gwern.net/doc/cs/security/2024-johnson.pdf}.
{\small\textit{Earliest verified question/agenda reference; historical priority qualified below. Requests rigorous privacy-benefit measurement and study of effective regulation and unintended effects.}}
\end{itemize}
\end{document}
